\chapter{Appendix}

\section{Core HLSL Path Tracing Implementation}
\label{apx:pathtracer_code}

The following code details the exact DirectX Raytracing implementation used within \textit{ChronoLux} to evaluate the dosimetric simulation. The integration is explicitly split into direct sun evaluation and a secondary stochastic multi-sampling loop for indirect ambient bouncing.

\begin{lstlisting}[language=C++, caption={IrradianceBake.raytrace: The primary ray generation kernel mapping Target Mesh UVs to stochastic integration over the hemisphere.}, label={lst:raygen}]
[shader("raygeneration")]
void IrradianceBakeRayGen()
{
    uint2 id = DispatchRaysIndex().xy;
    uint2 dims = DispatchRaysDimensions().xy;
    float4 bakedPos = _PositionMap[id];
    float4 bakedNormal = _NormalMap[id];
    
    if (bakedPos.w < 0.01 || length(bakedNormal.xyz) < 0.1) return;
    
    float3 origin = bakedPos.xyz;
    float3 normal = normalize(bakedNormal.xyz);
    
    uint baseSeed = Hash(id.y * dims.x + id.x + _FrameIndex);
    float totalIrradiance = 0.0;
    float nDotSun = dot(normal, _SunDirection);

    // 1. Direct Sun
    float directIrradiance = 0.0;
    if (nDotSun > 0.0)
    {
        uint sunSeed = baseSeed;
        directIrradiance = TracePath(origin, _SunDirection, normal, sunSeed, true, _SceneRTAS, _SunDirection, _BeamLux, _DiffuseLux) * nDotSun;
        totalIrradiance += directIrradiance;
    }

    // 2. Stochastic Multi-Sampling
    float indirectIrradiance = 0.0;
    uint numSamples = max(1, _SamplesPerPixel);
    for (uint i = 0; i < numSamples; i++)
    {
        uint sampleSeed = Hash(baseSeed + i);
        float3 randomDir = GetCosineWeightedDirection(normal, sampleSeed, PI);
        indirectIrradiance += TracePath(origin, randomDir, normal, sampleSeed, false, _SceneRTAS, _SunDirection, _BeamLux, _DiffuseLux);
    }
    
    float avgIndirect = (indirectIrradiance / (float)numSamples) * PI;
    totalIrradiance += avgIndirect;

    _DoseMap[id] += (totalIrradiance * _DeltaHours);
    _DirectDoseMap[id] += (directIrradiance * _DeltaHours);
    _IndirectDoseMap[id] += (avgIndirect * _DeltaHours);
}
\end{lstlisting}

\begin{lstlisting}[language=C++, caption={PathTracer.hlsl: The core recursive path tracing loop evaluating Russian Roulette absorption and transparent geometric traversal.}, label={lst:pathtracer}]
float TracePath(float3 origin, float3 direction, float3 surfaceNormal, inout uint seed, bool isDirectSun, RaytracingAccelerationStructure rtas, float3 sunDir, float beamLux, float diffuseLux)
{
    float throughput = 1.0;
    float3 currentOrigin = origin + surfaceNormal * RAY_BIAS;
    float3 currentDir = direction;

    [loop]
    for (uint bounce = 0; bounce < MAX_BOUNCES; bounce++)
    {
        RayDesc ray;
        ray.Origin = currentOrigin;
        ray.Direction = currentDir;
        ray.TMin = 0.001;
        ray.TMax = 10000.0;

        Payload payload;
        payload.isOccluded = 1;
        payload.reflectance = 0.0;
        payload.transmittance = 0.0;
        payload.hitDistance = 0.0;
        payload.worldNormal = float3(0, 1, 0);

        uint rayFlags = RAY_FLAG_FORCE_OPAQUE;
        if (isDirectSun) rayFlags |= RAY_FLAG_CULL_BACK_FACING_TRIANGLES;
        
        TraceRay(rtas, rayFlags, 0xFF, 0, 1, 0, ray, payload);

        if (payload.isOccluded == 0)
        {
            return EvaluateSky(currentDir, sunDir, beamLux, diffuseLux, isDirectSun) * throughput;
        }

        if (isDirectSun)
        {
            float trans = saturate(payload.transmittance);
            if (trans <= 0.0) return 0.0; // Blocked by opaque geometry
            
            throughput *= trans;
            // Use a larger bias to guarantee escaping the pane geometry, especially at grazing angles
            currentOrigin += currentDir * (payload.hitDistance + RAY_BIAS * 5.0);
            continue;
        }

        // Probabilistic event selection (Russian Roulette)
        float rv = Random(seed);
        float refl = saturate(payload.reflectance);
        float trans = saturate(payload.transmittance);
        
        if (rv < refl)
        {
            float3 hitPoint = currentOrigin + currentDir * payload.hitDistance;
            currentDir = GetCosineWeightedDirection(payload.worldNormal, seed, PI);
            currentOrigin = hitPoint + payload.worldNormal * RAY_BIAS;
        }
        else if (rv < refl + trans)
        {
            currentOrigin += currentDir * (payload.hitDistance + RAY_BIAS);
        }
        else
        {
            return 0.0;
        }

        if (throughput <= MIN_VISIBILITY) break;
    }
    return 0.0;
}
\end{lstlisting}

\section{C\# Orchestrator Implementations}
\label{apx:csharp_code}

The following code snippets detail the core C\# CPU orchestration logic. Listing \ref{lst:timeloop} demonstrates the dynamic chronological time-stepping Coroutine used to drive the simulation safely across the OS thread schedules.

\begin{lstlisting}[language={[Sharp]C}, caption={LightDoseSimulator.cs: The Coroutine driving dynamic time-stepping, asynchronous telemetry extraction, and synchronous EXR flushing.}, label={lst:timeloop}]
private IEnumerator RunSimulationInternal() {
    float deltaHours = stepSeconds / 3600f;
    int totalDays = endDay - startDay + 1;
    int currentDayIdx = 0;

    string exportDir = Path.Combine(Application.dataPath, "../Exports", DateTime.Now.ToString("yyyyMMdd_HHmmss"));
    Directory.CreateDirectory(exportDir);

    for (int day = startDay; day <= endDay; day++) {
        DateTime date = new DateTime(year, 1, 1).AddDays(day - 1);
        currentDayIdx++;
        
        DateTime sunrise = SunCalculator.FindSunrise(date, latitude, longitude, utcOffset);
        DateTime sunset = SunCalculator.FindSunset(date, latitude, longitude, utcOffset);
        
        if (sunrise == date && sunset == date) continue; // Polar night fallback
        
        for (DateTime localTime = sunrise; localTime < sunset; localTime = localTime.AddSeconds(stepSeconds)) {
            if (!_isSimulating) yield break;
            
            ApplySunPosition(localTime);
            
            // Optimization: Skip exact sunrise/sunset steps where there is functionally no light
            if (CurrentBeamLux < 0.1f && CurrentDiffuseLux < 0.1f) {
                continue;
            }

            // Dispatch hardware raytracing for this specific hour
            irradianceBaker.DispatchRays(CurrentSunDirection, CurrentBeamLux, CurrentDiffuseLux, deltaHours, samplesPerPixel, baker.PositionMap, baker.NormalMap);
            completedSteps++;
            currentProgress = ((float)currentDayIdx / totalDays) * 100f;
            
            // Extract telemetry data for the CSV asynchronously
            TakeAsyncSnapshot(day, localTime.ToString("yyyy-MM-dd HH:mm"), currentAltitude, currentAzimuth, CurrentBeamLux, CurrentDiffuseLux);

            if (completedSteps % 10 == 0) { 
                ApplyDosePreview(); 
            }
            // Yield execution back to the OS and Unity Main Thread to maintain interactivity
            yield return null;
        }
        
        // Flush the daily accumulation map to disk as a floating-point EXR
        ExportDailyMap(exportDir, day);
    }
}
\end{lstlisting}

Listing \ref{lst:objloader} showcases the background threaded parser utilized to load massively dense geometric heritage assets natively at runtime.

\begin{lstlisting}[language={[Sharp]C}, caption={RuntimeModelLoader.cs: Multi-threaded `.obj` chunk parsing avoiding main-thread freezes, followed by strict spatial normalization constraints.}, label={lst:objloader}]
string path = Path.Combine(ProjectManager.ModelFolder, artifactFile);
UpdateProgress(0.1f, $"Loading artifact: {artifactFile}");

// 1. Offload heavy string parsing to a background CPU thread
var parseTask = Task.Run(() => {
    if (!File.Exists(path)) return null;
    return SimpleObjLoader.ParseObj(File.ReadAllText(path));
});

// 2. Yield the main thread to keep the "LOADING ASSETS" UI animation smooth
while (!parseTask.IsCompleted) {
    UpdateProgress(0.1f, "Parsing artifact geometry...");
    yield return null;
}

if (parseTask.IsCompletedSuccessfully && parseTask.Result != null && parseTask.Result.Count > 0) {
    var dataList = parseTask.Result;
    _artifactInstance = new GameObject(Path.GetFileNameWithoutExtension(artifactFile));
    
    if (_artifactInstance != null) {
        _artifactInstance.transform.SetParent(modelRoot);
        _artifactInstance.transform.localPosition = Vector3.zero;
        _artifactInstance.transform.localScale = Vector3.one;

        // 3. Assemble the hierarchy on the main thread
        foreach (var data in dataList) {
            GameObject child = CreateGameObjectFromData(data, artifactMaterial);
            if (child == null) continue;
            child.transform.SetParent(_artifactInstance.transform, false);
            child.transform.localPosition = Vector3.zero;
            SetupArtifactPart(child, project); 
        }

        SetupArtifactRoot(_artifactInstance);
        
        // 4. Normalize the raw mesh geometry (Centering + Scaling children to a 2m box)
        CenterAndScaleModel(_artifactInstance, 2.0f);

        // 5. Apply user-defined transformations on the root
        _artifactInstance.transform.localPosition = project.artifactPosition;
        if (project.artifactScale != Vector3.zero) 
            _artifactInstance.transform.localScale = project.artifactScale;
    }
}
\end{lstlisting}